\chapter{Literature Review}

\section{Introduction}

This chapter critically reviews the research underpinning the project's
representation, compatibility-modelling, retrieval and evaluation choices.
Fashion recommendation includes visual similar-item retrieval, substitute
recommendation, personalised ranking, outfit compatibility prediction and
complementary-item retrieval \parencite{deldjoo2023review,sarkar2023outfittransformer}.
Because these tasks differ in their inputs, supervision and definitions of
success, evidence obtained for one cannot automatically support another.

The review first positions complete-outfit recommendation within recommender
systems, then examines the development of fashion-compatibility models and the
role of frozen vision--language representations. It subsequently considers
learning objectives, evaluation, candidate search and responsible design. The
literature is used to derive controlled alternatives, not to assume that the
newest or most complex architecture is necessarily appropriate.

\section{Positioning complete-outfit recommendation}

Traditional recommender systems often learn from interactions such as ratings,
clicks or purchases and ask which item should be presented to a particular
user. Evaluation must be connected to the user task, available observations and
intended use; otherwise, nominally similar experiments may not be comparable
\parencite{herlocker2004evaluation}. Fashion recommenders extend this setting
with visual, textual, contextual and social information
\parencite{deldjoo2023review}.

Cove has no longitudinal user--item history and therefore does not claim to
learn stable individual preferences or compare collaborative-filtering
algorithms. Its learned task is content-based outfit compatibility: estimating
whether relationships within a set of garments resemble those observed in the
training outfits. At runtime, that score ranks completions around a
user-selected garment. Dataset compatibility and personal preference must
remain distinct, because following a dataset regularity does not establish what
one person will prefer.

Complete-outfit recommendation is consequently considered at three connected
levels. Representation converts garment images into features; relational
scoring evaluates garments jointly; decision and interaction filter infeasible
candidates, search combinations and expose alternatives. Secondary outfit data
can evaluate the first two levels. The third can be tested for deterministic
behaviour, completeness and runtime, but not for satisfaction without human
evidence. This separation prevents the compatibility classifier from being
treated as the entire recommender system.

\section{Development of fashion-compatibility models}

Early image-based recommendation distinguished visual similarity from
complementarity. Items serving the same role may be substitutes, whereas items
from different roles may be complementary despite being visually dissimilar
\parencite{mcauley2015styles}. Pairwise metric learning operationalises this
distinction by learning spaces in which compatible pairs are close. Such models
support efficient nearest-neighbour retrieval, but a complete outfit contains
several interacting pairs. Summing local scores assumes that pairwise
compatibility is sufficient to represent the set.

\textcite{han2017bilstm} represented outfits as ordered garment sequences and
used a bidirectional long short-term memory network. This aggregates information
from garments on both sides of an item, although outfits have no natural
linguistic order and must be ordered through garment categories. The method
therefore provides a useful structured baseline while retaining possible
sensitivity to that imposed sequence.

\textcite{vasileva2018typeaware} argued that compatibility depends on item type
and proposed type-aware embedding spaces. A top--shoe relationship need not use
the same transformation as a top--top relationship, making category or slot
identity part of the learned relation rather than descriptive metadata. This
supports testing slot-pair-specific projections and explicit slot
representations.

Context-aware graph modelling provides another response to isolated pair
relations. \textcite{cucurull2019context} generated product embeddings
conditioned on known compatible context and reported stronger compatibility
prediction when more context was available. The finding supports the view that
compatibility is conditional rather than intrinsic to one item. Graph
construction requires available nodes and edges, which may not exist for a
newly uploaded garment.

Attention allows each garment representation to interact with every other
garment without requiring a fixed sequence. OutfitTransformer uses task-specific
tokens and self-attention for outfit compatibility and complementary-item
retrieval \parencite{sarkar2023outfittransformer}. Treating an outfit as an
unordered set is appropriate because rearranging the same garments should not
change compatibility. Attention offers flexible higher-order interaction but
also adds capacity. Results obtained under OutfitTransformer's original
protocol do not establish that equal capacity is required with frozen inputs
and a smaller processed dataset.

These studies support several modelling assumptions rather than one universal
architecture. Pairwise models prioritise local relations and efficiency;
recurrent models learn structured aggregation; type-aware models encode
category-dependent relations; and attention learns set-level interactions
\parencite{han2017bilstm,vasileva2018typeaware,sarkar2023outfittransformer}.
Their differences motivate controlled comparisons using the same features,
splits and training conditions.

\section{Polyvore data and label construction}

Polyvore datasets convert outfits assembled by platform users into observable
item sets and are consequently practical secondary data for compatibility
research \parencite{han2017bilstm,vasileva2018typeaware}. A positive example
records an outfit retained in the dataset; it is not evidence of an objective
or universally shared aesthetic judgement.

Negative outfits are commonly constructed by replacing an item in a positive
outfit. Same-type replacement is harder than arbitrary-category replacement
because the negative remains structurally plausible
\parencite{vasileva2018typeaware}. Since replacements are generated rather than
independently judged incompatible, they may include false negatives. This risk
follows from the published construction procedure; it does not establish that
a known proportion of labels is incorrect.

Splitting the data changes the generalisation question. If a product appears in
both training and testing outfits, a model may benefit from recognising an
observed item. Vasileva et al. describe their revised Polyvore partition as a
disjoint split intended to reduce this problem
\parencite{vasileva2018typeaware}. The present project retains that published
name, while independently auditing the locally acquired metadata rather than
assuming complete item-level separation. Such a split also leaves shared
photographic conventions, category distributions, platform culture and
labelling assumptions unchanged.

In the common fill-in-the-blank (FITB) task, one garment is removed and the
model selects the correct completion from four candidates
\parencite{han2017bilstm,vasileva2018typeaware}. FITB evaluates relative ranking
within a small set and is closer to recommendation than binary classification,
but remains smaller than full-catalogue retrieval. OutfitTransformer evaluates
compatibility prediction and large-scale retrieval separately, reinforcing
that FITB does not establish scalable retrieval quality
\parencite{sarkar2023outfittransformer}.

Polyvore is therefore suitable for controlled training and offline evaluation.
It cannot represent all clothing practices. Filtering official questions to
four supported slots introduces an additional coverage boundary. The retained
fraction must accompany the subset result so that it is not presented as the
complete official benchmark.

\section{Frozen vision--language representations}

Earlier compatibility systems commonly used convolutional encoders pretrained
on ImageNet and adapted to fashion data \parencite{vasileva2018typeaware}.
Vision--language models instead learn from large collections of image--text
pairs. CLIP demonstrated transferable visual representations and zero-shot
recognition across numerous downstream datasets \parencite{radford2021clip}.

SigLIP belongs to this family but replaces a global softmax contrastive
objective with independent sigmoid losses over image--text pairs
\parencite{zhai2023siglip}. It was not pretrained for outfit compatibility; its
image encoder instead supplies a general visual representation from which a
smaller downstream model can attempt to learn garment relationships.

Freezing the encoder converts each image into a fixed vector and restricts
training to the compatibility head. For this project, that reduces trainable
parameters, repeated image-processing cost and representation variation across
baselines. It also permits comparisons with identical inputs. These are
computational and experimental advantages within the project, not evidence
that frozen features are universally preferable.

The downstream head cannot recover information absent from the frozen vector.
Transfer is not guaranteed by source-model scale alone:
\textcite{kornblith2019transfer} found that transfer depends on both
representation and target task, with limited benefit on some fine-grained
datasets. Web-scale pretraining can also transmit biases in the source data
\parencite{radford2021clip,mehrabi2021bias}. Frozen SigLIP is thus a
resource-conscious hypothesis to test rather than an assumed fashion-specific
solution.

\section{Learning objectives and controlled baselines}

Compatibility classification and recommendation ranking are related but
distinct. Binary cross-entropy trains separation of labelled positive and
negative outfits. Its score can order candidates, but the loss does not
directly optimise order among competing completions. Pairwise ranking instead
encourages a preferred observation to score above an alternative; Bayesian
Personalised Ranking is one established formulation for implicit observations
\parencite{rendle2009bpr}.

Neither objective is automatically preferable here. Binary labels are directly
available in the compatibility rows, whereas ranking requires defensible
positive--negative pairs. If a synthetic negative is noisy, a ranking objective
may optimise an unreliable comparison. The choice must be evaluated rather
than inferred from the fact that runtime recommendation uses ranking.

The controlled baselines form a hierarchy. Chance provides a lower bound for
four-choice FITB. Non-learned cosine similarity tests whether frozen features
already contain useful structure. A linear model over pooled features tests
whether learned weighting is sufficient without nonlinear interaction. These
controls prevent the downstream architecture from receiving credit for
information already present in SigLIP.

A recurrent baseline represents the sequence assumption of
\textcite{han2017bilstm}; a type-aware model represents relation-specific
transformations from \textcite{vasileva2018typeaware}; and attention represents
set-level interaction following \textcite{sarkar2023outfittransformer}. These
are controlled approximations, not exact reproductions, which would require the
original encoders, preprocessing, objectives and hyperparameters.

Structural ablation addresses component contribution. Removing the pairwise
module tests whether explicit item--item interaction adds evidence beyond
pooled outfit information. Removing slot embeddings tests whether slot identity
adds evidence beyond visual vectors. If removal does not produce a consistent
decline across repeated runs, the component's benefit remains uncertain.

\section{Evaluation beyond a single score}

Evaluation should match the task claimed \parencite{herlocker2004evaluation}.
Compatibility classification, candidate completion, confidence interpretation
and application behaviour therefore require separate evidence.

ROC--AUC measures discrimination across thresholds and can be interpreted as
the probability that a randomly selected positive receives a higher score than
a randomly selected negative \parencite{fawcett2006roc}. It does not measure
whether the correct completion ranks first in a candidate set, nor whether a
score is a calibrated probability.

FITB evaluates completion more directly. In a balanced four-choice question,
chance accuracy is 25\%, but performance still depends on candidate and label
construction. Correct answers must be matched through the documented dataset
relationship rather than assumed to occupy a fixed array position. Retaining
question-level outputs makes label correction and paired comparison auditable.

Calibration is separate from discrimination. Neural networks may rank examples
correctly while remaining overconfident; temperature scaling was effective in
several classification settings examined by \textcite{guo2017calibration}.
Because a positive scalar temperature does not alter logit order, it can improve
confidence interpretation without improving FITB ranking.

Stochastic training variation must also be reported. Parameter initialisation,
data order and other benchmark stages can alter an otherwise unchanged
experiment \parencite{bouthillier2021variance}. Cove uses three seeded runs for
its principal architecture comparisons. This is a resource-constrained
repeated-run design, not a complete characterisation of training variance.

Structural ablations are evaluated on the same FITB questions and matched
seeds. Pairing preserves the full--ablation relationship on each evaluation
unit instead of treating their predictions as independent
\parencite{dietterich1998tests}. The project's bootstrap averages the three
seeded outcomes within each question and then resamples questions, following
the bootstrap principle of estimating sampling uncertainty by resampling
\parencite{efron1994bootstrap}. Its interval is conditional on those runs and
does not represent every possible training seed.

Overfitting is examined through training, validation and held-out test
behaviour. Validation-based early stopping is intended to stop iterative
training when further training improvement is no longer accompanied by better
validation performance \parencite{prechelt1998early}. A small
validation--test difference supports only the consistency of the selected
checkpoint on the current test split, not general reliability. Stronger
training than validation and test performance is consistent with training-set
overfitting, although data composition can also contribute.

No offline metric establishes user benefit. Without participant evidence,
claims remain limited to discrimination, controlled completion ranking,
calibration, runtime, coverage and deterministic system behaviour.

\section{Candidate search and explicit constraints}

A model scores only candidates that reach it. If a strong candidate is removed
by a frequency quota or similarity shortlist, later ranking cannot recover it.
Candidate generation is therefore part of recommendation behaviour rather than
neutral preprocessing.

Large-scale complementary-item systems commonly use embedding retrieval.
OutfitTransformer creates a target representation and applies nearest-neighbour
retrieval \parencite{sarkar2023outfittransformer}. Vector indexes and
approximate search make this efficient at very large scale, trading retrieval
completeness for speed or memory \parencite{johnson2021similarity}.

Cove instead assigns a non-additive score to a complete combination. An item
that is not individually closest to the uploaded garment can still participate
in the highest-scoring complete outfit, so independent slot shortlists may not
preserve the final optimum. This follows from the deployed scoring function
rather than from a claim made by the retrieval literature.

For a bounded four-slot space, batched exhaustive scoring can remain feasible.
Here, \emph{exact} means that every combination surviving the declared filters
is evaluated under one deployed objective. A safety limit is required so an
oversized search fails explicitly rather than presenting a partial prefix as
complete.

Weather suitability, slot occupancy and the user-selected clothing range are
feasibility requirements, not aesthetic preferences inferred from Polyvore.
The model ranks feasible outfits, while deterministic rules define what is
allowed. This keeps model behaviour, search completeness and functional
constraints separately auditable.

\section{Bias, personalisation and responsible design}

Bias can enter through data selection, labels, pretrained representations,
candidate construction and interaction design, and cannot be reduced to
overall accuracy \parencite{mehrabi2021bias}. Clothing practice is also
cultural, temporal and subjective \parencite{deldjoo2023review}.

Polyvore reflects a historical platform population and may over-represent some
categories, presentation styles and commercial photography. Synthetic
same-type negatives add another assumption, while SigLIP carries uncertainty
from web-scale pretraining. These sources justify an audit and explicit
limitations, but do not permit demographic fairness to be measured without
demographic attributes.

Automatic gender inference would add an uncertain identity classification, a
practice criticised in human--computer interaction research
\parencite{keyes2018misgendering}. Cove asks the user to choose a menswear or
womenswear candidate range instead. Those terms remain simplified catalogue
filters and are not definitions of the user.

Local execution reduces external transmission of wardrobe data but does not
resolve feedback bias. Exposure influences which clicks become observable, and
simulation evidence shows that algorithmically confounded feedback can
increase behavioural homogeneity without increasing utility
\parencite{chaney2018confounding}. Defensible long-term self-learning would
need exposure logging, exploration, feedback-loop controls and separate
evaluation, so it lies outside the frozen study.

Fairness cannot be detached from social context or treated solely as a property
of an isolated component \parencite{selbst2019fairness}. A proportionate
response documents limitations, avoids identity inference, retains the
selected garment, exposes alternatives and permits rejection or replacement.

\section{Synthesis and research framework}

The literature supports four conclusions. First, compatibility is relational:
visual similarity is insufficient and category identity can form part of the
relation, although maximum model complexity is not always justified
\parencite{mcauley2015styles,vasileva2018typeaware,sarkar2023outfittransformer}.
This supports evaluating a lightweight relational head against progressively
stronger alternatives.

Second, frozen vision--language features offer a feasible representation
strategy rather than a fashion-specific solution. They reduce training burden
while retaining transfer, domain and inherited-bias risks
\parencite{radford2021clip,zhai2023siglip,kornblith2019transfer}. The downstream
head must consequently be compared with non-learned and learned controls.

Third, ROC--AUC, FITB, calibration, repeated runs, ablation, coverage and domain
diagnostics answer different questions
\parencite{fawcett2006roc,guo2017calibration,bouthillier2021variance}.

Fourth, candidate generation and explicit constraints affect deployed
behaviour as well as scoring. Approximate retrieval is appropriate at very
large scale, but a bounded non-additive four-slot objective creates a testable
case for exact constrained search
\parencite{sarkar2023outfittransformer,johnson2021similarity}.

The remaining empirical questions are whether a small relational head on
frozen SigLIP features can balance discrimination, completion ranking, runtime
and auditability under one secondary-data protocol, and how much deployed-model
quality is lost through quota candidate generation. Chapter 3 translates those
questions into controlled experiments.

\section{Summary}

Cove has been positioned as a content-based compatibility and constrained-search
system rather than a conventional personalised collaborative recommender. The
literature shows a progression from pairwise visual relations to sequence,
type-aware, contextual and attention-based modelling. It supports frozen
vision--language features as a resource-conscious experimental layer while
requiring explicit treatment of transfer, labels, evaluation and bias.

Chapter 3 applies this framework through the methodology used for model and
system evaluation.
